%************************************************
\chapter{Old Laws, New Results}\label{ch:alignment}
%************************************************

\chapterlead{A new theory owes two debts. It must say again, exactly, what the old laws say wherever they have been tested, and it must then say something they do not, which a measurement could prove wrong. This chapter states the two debts, shows how much of the first this volume pays, and names the candidate for the second.}

\scene{London, 6 November 1919}

On 6 November 1919 the rooms of the Royal Society in London were full. The Society and the Royal Astronomical Society had called a joint meeting to hear the results of two expeditions sent to watch the total eclipse of the Sun of 29 May, one to Sobral, in the north of Brazil, the other to the island of Pr{\'i}ncipe, in the Gulf of Guinea. Each team had photographed the stars around the darkened Sun and compared their places with photographs of the same stars taken when the Sun was far from them \citep{Dyson1920}.

Starlight that grazes the Sun should bend. Newton's law of gravity, applied to light as if it were a falling body, gave a bend of $0.87$ seconds of arc at the edge of the Sun. Einstein's new theory of gravity gave twice as much, $1.75$ seconds \citep{Einstein1916,Dyson1920}. Frank Dyson, the Astronomer Royal, gave the result of the Sobral plates: $1.98$ seconds, with a probable error of $0.12$. Arthur Eddington, back from Pr{\'i}ncipe, gave his: $1.61$ seconds, with a probable error of $0.30$. A third set of plates, from a second telescope at Sobral, gave a smaller value; the team judged its star images faulty and set it aside, a decision that was argued over for years. Later measurements with radio telescopes have confirmed Einstein's value to better than one part in ten thousand \citep{Will2014}.

J.~J. Thomson, who presided, called it \enquote{the most important result obtained in connection with the theory of gravitation since Newton's day} \citep{Observatory1919}. The philosopher Alfred North Whitehead was in the room. \enquote{The whole atmosphere of tense interest was exactly that of the Greek drama,} he wrote six years later, and he remembered, in the background, \enquote{the picture of Newton to remind us that the greatest of scientific generalisations was now, after more than two centuries, to receive its first modification} \citep{Whitehead1925}.

Two things made that meeting decisive, and their order matters. First, Einstein's theory gave back Newton's. For slow bodies in weak gravity its equations become Newton's law, so it agreed with two centuries of measurements of falling bodies, tides, planets and comets \citep{Einstein1916}. It also settled an old puzzle: the orbit of Mercury turns by $43$ seconds of arc per century more than Newton's law allows, even with the pull of all the other planets, and in 1915 Einstein's theory gave exactly that \citep{Einstein1915}. But the turning of Mercury was known before the theory, which explained it after the fact. Second, the theory said something new, before anyone had measured it: light grazing the Sun bends by $1.75$ seconds, not by $0.87$. That number could have come out wrong. It did not.

\bigskip

This chapter is about those two duties of a new theory, and about where this volume stands with each. \Cref{sec:al-duties} states them, with five theories that met both. \Cref{sec:al-rules} sets out the rules under which this volume tests itself. \Cref{sec:al-found} gathers what the cosmic mass language has been found to say, the laws of Newton, Coulomb, Maxwell and Lorentz, Heaviside's gravity, the balance of nuclear fusion, and a mass first written down by George Johnstone Stoney, and asks what it means for physics to have all of them in one language. \Cref{sec:al-coupling} turns to the Coupling Force, the candidate for new results, which is the subject of Volume~2.

%------------------------------------------------
\section{The two duties of a new theory}\label{sec:al-duties}
%------------------------------------------------

The first duty is to agree. A new theory must give back every old law wherever the old law has been tested, with the same measured constants. Niels Bohr made this a rule for the quantum theory of the atom, which had to give back the classical laws of radiation for slow vibrations and large orbits; in 1920 he gave it a name, the principle of correspondence \citep{Bohr1920}. The great theories since have kept it. Relativity gives back Newton's mechanics for slow bodies and Newton's gravity for weak fields; quantum mechanics gives back classical mechanics for large bodies. The old law is not thrown away. It becomes a special case, true inside a domain whose borders the new theory draws.

The second duty is to risk being wrong. Karl Popper put it in one sentence: \enquote{it must be possible for an empirical scientific system to be refuted by experience} \citep[section~6]{Popper1959}. A theory that agrees with everything known but forbids nothing has told us nothing. Imre Lakatos added a measure of progress for a whole line of research: it advances when its new versions predict new facts, and some of those facts are then found; it degenerates when it only adjusts itself, after the event, to facts that others found \citep{Lakatos1970}.

% keep the statement of the two duties together: break the page here if fewer than 7 lines are left
\par\vskip 0pt plus 7\baselineskip\penalty-100\vskip 0pt plus -7\baselineskip\vskip 7\baselineskip\penalty9999\vskip -7\baselineskip\vskip 0pt
Put together, the two duties read:
\begin{quote}
\itshape A new theory must give back every old law wherever that law has been tested, with the same constants; and it must predict at least one new result, a number with its error, stated before it is measured, that a measurement could prove wrong.
\end{quote}

\Cref{tab:al-history} shows five theories that met both duties. In each of them the old laws were kept whole, and the new result was not put in on purpose.

\begin{table}[tb]
\myfloatalign
\small
\begin{tabularx}{\textwidth}{@{}>{\raggedright\arraybackslash}p{0.2\textwidth} >{\raggedright\arraybackslash}X >{\raggedright\arraybackslash}X >{\raggedright\arraybackslash}X@{}}
\toprule
\tableheadline{Theory} & \tableheadline{Old laws kept} & \tableheadline{New result} & \tableheadline{Found} \\
\midrule
Newton's gravity, in Le Verrier's hands, 1846 & the motions of the known planets & an unseen planet beyond Uranus, and where to look for it & Neptune, by Galle and d'Arrest, 23 September 1846, within one degree of the place \citep{LeVerrier1846} \\
Maxwell's equations, 1865 & the laws of Coulomb, Amp{\`e}re and Faraday & electric waves that travel at the speed of light & made and detected by Hertz, 1887 to 1888 \citep{Maxwell1865,Hertz1893} \\
Mendeleev's periodic law, 1871 & the known elements and their chemistry & an element below aluminium, with its weight and properties & gallium, by Lecoq de Boisbaudran, 1875 \citep{Mendeleev1871,Lecoq1875} \\
Einstein's gravity, 1915 & Newton's law for slow bodies and weak fields & starlight bent by $1.75$ seconds of arc at the edge of the Sun & the eclipse of 1919 \citep{Einstein1916,Dyson1920} \\
Dirac's electron, 1928 to 1931 & quantum mechanics and special relativity & a particle with the mass of the electron and the opposite charge & the positron, by Anderson, 1932 \citep{Dirac1928,Dirac1931,Anderson1933} \\
\bottomrule
\end{tabularx}
\caption[Five theories that met both duties]{Five theories that met both duties. Each kept the old laws whole, and each predicted a result that nobody had seen and that a measurement could have refuted.}
\label{tab:al-history}
\end{table}

Every row has the same shape. The new theory first contained the old one, and only then did its new prediction mean something, because it came from a theory that was already right wherever it had been checked. A theory that only agrees is a new way of writing known physics: a language. A theory that also predicts, and survives the test, is new physics.

\begin{plainwords}
A new theory has two jobs. First, it must get the old answers right: wherever the old laws were measured and found true, the new theory must give the same numbers. Second, it must say something new that could turn out to be false, and say it before anyone measures it. Einstein's theory did both. It gave Newton's answers for the planets, and it said that starlight would bend twice as much as Newton's law said. The eclipse of 1919 could have proved it wrong, and it did not.
\end{plainwords}

%------------------------------------------------
\section{The rules this volume keeps}\label{sec:al-rules}
%------------------------------------------------

This volume holds itself to the first duty, under rules that it fixes in advance.
\begin{enumerate}
  \item \emph{Two inputs.} Each body enters with two measured numbers, its mass in kilograms and its charge in coulombs, and with nothing else (\Cref{tab:scope}).
  \item \emph{Measured constants only.} The constants are those of physics: $G$, $K$, $c$, $e$ and $\hbar$ (\Cref{app:notation}). The size of the Composer Unit, $|\Gp| = \sqrt{K/G}$, is not fitted to anything; $G$ and $K$ fix it (\Cref{subsec:cu-size}).
  \item \emph{One reading.} Every measured force, energy and momentum is read as the real part of a General Number. The rule is stated once, in Chapter~\ref{ch:interaction} (Definition~\ref{def:reading}), and used without exception afterwards.
  \item \emph{One angle.} The angle of the Composer axis is held at $79.86^\circ$ from Chapter~\ref{ch:interaction} to the end of the book (\Cref{subsec:which-root}). It is a choice of coherence, not a fit to data.
  \item \emph{Every number checked.} Each result is set beside a measurement, or beside the value of the standard laws, with its source, and the difference is reported.
\end{enumerate}
Under these rules, each agreement is a test that could have failed. The last rule also says what an agreement is worth. When the language gives back a known law, it has passed a test of alignment, and it has discovered nothing yet. Chapter~\ref{ch:unified} says so of its own results: each is a check on the reading rather than a discovery.

\begin{table}[p]
\myfloatalign
\footnotesize
\begin{tabularx}{\textwidth}{@{}>{\raggedright\arraybackslash}p{0.26\textwidth} >{\raggedright\arraybackslash}X >{\raggedright\arraybackslash}p{0.29\textwidth}@{}}
\toprule
\tableheadline{Old law} & \tableheadline{In cosmic masses} & \tableheadline{Where, how well} \\
\midrule
\multicolumn{3}{@{}l}{\emph{Bodies at rest, and matter}} \\[2pt]
the gravity of \citet{Newton1687} and the law of Coulomb~\citeyearpar{Coulomb1785} & the real part of one product of two cosmic masses & Proposition~\ref{prop:law-real}: both exactly, each with its full constant, at every angle \\
the rule of \citet{DuFay1733}: like charges push apart & the minus sign of the Rule of the Square & Proposition~\ref{prop:square-rule}: in every true plane \\
the hydrogen atom of \citet{Bohr1913} & Coulomb's law at its full strength & \Cref{subsec:test-sign}: $13.598$ eV, as measured \citep{Sansonetti2005} \\
the neutral neutron & the Coupling term has no real part & \Cref{subsec:test-sign}: zero; measured, $(-0.2 \pm 0.8) \times 10^{-21}\,e$ \citep{Baumann1988} \\
the masses of \citet{Stoney1881} and \citet{Planck1899} & the charge part of one elementary charge, $m_S = e\,|\Gp|$ & Definition~\ref{def:stoney-mass}, Remark~\ref{rem:stoney-planck}: $m_S = \sqrt{\alpha_{\mathrm{fs}}}\,m_P$ \\
the charge-to-mass ratios of \citet{Thomson1897} and \citet{Aston1919} & one angle for each body, on one frame & Chapter~\ref{ch:standard-model}: $306$ bodies, no measured value moved \\
charge and mass in fusion; the mass relations of \citet{GarveyKelson1966} & two coordinates of one conservation law; six neighbours whose charge parts cancel exactly & Chapter~\ref{ch:closure}: the masses of $2{,}481$ nuclei to $0.25$ MeV, and to $0.11$ MeV for $A \ge 40$ \\
\midrule
\multicolumn{3}{@{}l}{\emph{Bodies in motion, and fields}} \\[2pt]
the equations of \citet{Maxwell1865} & the Composer coordinate of the field equations, with $\mu_0 = 4\pi K/c^2$ & \Cref{thm:un-maxwell}: exactly \\
Newton's mechanics, and the gravity of \citet{Heaviside1893} & the real coordinate of the same equations, and of the law of motion & Law~\ref{law:unified}: as Newton and Heaviside \\
the force of Lorentz & the real part of $\mu\,(\boldsymbol{\mathcal{E}} + \mathbf{v}\times\boldsymbol{\mathcal{B}})$ & Proposition~\ref{prop:un-lorentz}: Thomson's tube, the cyclotron, the drift in crossed fields \\
waves at the speed of light & both coordinates of the fields & \Cref{thm:un-waves}: gravity as fast as light within $3$ parts in $10^{15}$ \citep{Abbott2017GRB} \\
the conservation of energy and of momentum & one balance for both parts; one action for everything & \Cref{thm:un-energy,thm:un-action} \\
the Moon, and hydrogen in motion & one program for every body & \Cref{tab:un-results}: $27.28$ days against $27.32$; $13.598$ eV \\
\bottomrule
\end{tabularx}
\caption[The old laws in the cosmic mass language]{The old laws in the cosmic mass language: where each is found in this volume, and how closely it agrees.}
\label{tab:al-found}
\end{table}

%------------------------------------------------
\section{What the cosmic mass language already says}\label{sec:al-found}
%------------------------------------------------

Chapters~\ref{ch:interaction} to~\ref{ch:unified} test the language against the laws of physics, one after another. \Cref{tab:al-found} gathers what they find, and \Cref{fig:al-map} draws it as a map: one number for each body, one product for each pair, one way of reading it, and inside them the old laws. This section walks through the map in four steps, from two bodies at rest to moving bodies and their fields, and then asks what the result is worth.

\begin{figure}[tb]
\myfloatalign
\begin{tikzpicture}[>=Stealth,
    law/.style={draw=RoyalBlue!55, fill=RoyalBlue!6, rounded corners=2pt, text width=3.1cm,
                align=flush center, inner sep=3pt, font=\footnotesize, minimum height=1.2cm},
    core/.style={draw=RoyalBlue!80, fill=RoyalBlue!14, rounded corners=3pt, text width=3.45cm,
                align=flush center, inner sep=4pt, font=\footnotesize, minimum height=1.25cm},
    open/.style={draw=Maroon!60, densely dashed, fill=Maroon!4, rounded corners=2pt, text width=3.45cm,
                align=flush center, inner sep=3pt, font=\footnotesize, minimum height=1.2cm},
    tag/.style={font=\scriptsize, text=gray!70!black},
    link/.style={->, thick, RoyalBlue!60}]
  % the one law, in the middle
  \node[core] (c) at (5.9,0) {\textbf{one number per body}\\ one product per pair,\\ read by its real part};
  % bodies at rest, on the left
  \node[tag] at (1.7,3.2) {\textsc{bodies at rest}};
  \node[law] (l1) at (1.7,2.25) {Newton's gravity and Coulomb's law\\ {\color{gray!75!black}Chapter~\ref{ch:interaction}}};
  \node[law] (l2) at (1.7,0.75) {like charges push apart\\ {\color{gray!75!black}Chapter~\ref{ch:interaction}}};
  \node[law] (l3) at (1.7,-0.75) {hydrogen; the neutral neutron\\ {\color{gray!75!black}Chapter~\ref{ch:interaction}}};
  \node[law] (l4) at (1.7,-2.25) {Stoney's mass; $306$ bodies\\ {\color{gray!75!black}Chapter~\ref{ch:standard-model}}};
  % bodies in motion, on the right
  \node[tag] at (10.1,3.2) {\textsc{bodies in motion}};
  \node[law] (r1) at (10.1,2.25) {Maxwell's equations\\ {\color{gray!75!black}Chapter~\ref{ch:unified}}};
  \node[law] (r2) at (10.1,0.75) {Newton's mechanics; Heaviside's gravity\\ {\color{gray!75!black}Chapter~\ref{ch:unified}}};
  \node[law] (r3) at (10.1,-0.75) {the force of Lorentz\\ {\color{gray!75!black}Chapter~\ref{ch:unified}}};
  \node[law] (r4) at (10.1,-2.25) {waves at the speed of light; energy\\ {\color{gray!75!black}Chapter~\ref{ch:unified}}};
  % nuclei, below
  \node[law] (n1) at (5.9,-2.25) {nuclei: charge and mass kept in fusion\\ {\color{gray!75!black}Chapter~\ref{ch:closure}}};
  % what the real part leaves, above
  \node[open] (v2) at (5.9,2.25) {\textbf{Composer coordinate}\\ Coupling term, angle\\ {\color{Maroon!80!black}Volume~2}};
  \foreach \t in {l1,l2,l3,l4} \draw[link] (c.west) -- (\t.east);
  \foreach \t in {r1,r2,r3,r4} \draw[link] (c.east) -- (\t.west);
  \draw[link] (c.south) -- (n1.north);
  \draw[->, thick, Maroon!60, densely dashed] (c.north) -- (v2.south);
\end{tikzpicture}
\caption[The old laws inside one law]{The old laws inside one law. Read by its real part, the product of two cosmic masses gives back the laws of bodies at rest (left), of bodies in motion (right) and of nuclei (below), each in the chapter shown. What the real part leaves, the Composer coordinate, holds the Coupling term and the angle of the Composer axis: the questions of Volume~2 (above).}
\label{fig:al-map}
\end{figure}

\subsection{One product, every sign}\label{subsec:al-sign}

Physics has always written gravity and electricity as two laws with two constants, and it adds a rule of signs by hand. Masses always pull. Charges pull when their signs differ and push when they agree, as Du Fay found in 1733 \citep{DuFay1733}. Nothing in the two laws says why: the rule lives in the signs given to the constants and to the charges, and it is learnt as a fact of nature.

In the cosmic mass language one law serves both, and no rule of signs is added. For two bodies at rest, the force is the real part of the product of their cosmic masses, multiplied by $G$ and divided by the square of the distance, and it is counted positive for a pull (Proposition~\ref{prop:law-real}):
\[
F = \frac{G\,\Re(\mu_a\mu_b)}{r^2} = \frac{G\,m_a m_b - K\,q_a q_b}{r^2} .
\]
The minus sign between the two terms is not a convention. It is the real part of the square of the Composer Unit, $\Re(\Gp^2) = -|\Gp|^2$, the same at every angle of a true plane (Proposition~\ref{prop:square-rule}): a unit that leans out of the line has a square that points backward. The pull between masses and the push between like charges come from the geometry of one plane.

\begin{workedexample}[Worked example: one product, three answers]
\begin{itemize}
  \item \emph{Two electrons, one metre apart.} The charge part of each is $e\,|\Gp| = m_S$, the Stoney mass, so $\Re(\mu_e\mu_e) = m_e^2 - m_S^2 \approx 8.3 \times 10^{-61} - 3.457 \times 10^{-18} \approx -3.457 \times 10^{-18}$ kg$^2$. Multiplied by $G$, the force is $-2.307 \times 10^{-28}$ N: a push, Coulomb's force between two electrons. The share of gravity in it is $2.4 \times 10^{-43}$.
  \item \emph{An electron and a proton, one metre apart.} Now $q_a q_b = -e^2$, and $\Re(\mu_e\mu_p) = m_e m_p + m_S^2 \approx 3.457 \times 10^{-18}$ kg$^2$: a pull of $2.307 \times 10^{-28}$ N.
  \item \emph{An apple of $0.1$ kg at the surface of the Earth.} Both bodies are neutral, so the real part is the product of the masses, and the force is $G \times 0.1\ \mathrm{kg} \times 5.972 \times 10^{24}\ \mathrm{kg} / (6.371 \times 10^{6}\ \mathrm{m})^2 \approx 0.98$ N: a pull, the weight of the apple.
\end{itemize}
One expression gave a push, a pull and a weight. Each time the sign came from the plane, and the size from the two parts of the cosmic masses.
\end{workedexample}

The same expression shows why electricity is so much stronger than gravity between particles. Between two electrons the push and the pull stand in the ratio $(m_S/m_e)^2 \approx 4.2 \times 10^{42}$, because the charge part of an electron is about $2 \times 10^{21}$ times its mass: its cosmic mass lies almost along the Composer axis. Physics has long quoted this ratio as a bare number. Here it is the same number, with a picture: the lean of one arrow.

\subsection{One number for every body}\label{subsec:al-bodies}

The Stoney mass of the example above shows the kind of coherence that the language brings. In 1874 Stoney proposed units of nature in which nothing would be arbitrary, built from three constants: the speed of light, the constant of gravity, and the portion of electricity that every broken chemical bond carries through an electrolyte, which he later named the electron \citep{Stoney1881,Stoney1894}. His unit of mass reappears here with a new meaning: it is the charge part of one elementary charge.

% keep the Stoney example on one page: break the page here if fewer than 15 lines are left
\par\vskip 0pt plus 15\baselineskip\penalty-100\vskip 0pt plus -15\baselineskip\vskip 15\baselineskip\penalty9999\vskip -15\baselineskip\vskip 0pt
\begin{workedexample}[Worked example: the Stoney mass]
\begin{itemize}
  \item \emph{The size of the Composer Unit.} $|\Gp| = \sqrt{K/G} \approx 1.16043 \times 10^{10}$ kg per coulomb (\Cref{subsec:cu-size}).
  \item \emph{The charge part of one elementary charge.} $m_S = e\,|\Gp| \approx 1.602177 \times 10^{-19}\ \mathrm{C} \times 1.16043 \times 10^{10}\ \mathrm{kg/C} \approx 1.85921 \times 10^{-9}$ kg (Definition~\ref{def:stoney-mass}).
  \item \emph{A balanced pair.} $G\,m_S^2 \approx 6.67430 \times 10^{-11} \times (1.85921 \times 10^{-9})^2 \approx 2.30708 \times 10^{-28}$ N\,m$^2$, and $K e^2 \approx 8.98755 \times 10^{9} \times (1.602177 \times 10^{-19})^2 \approx 2.30708 \times 10^{-28}$ N\,m$^2$. Two bodies of mass $m_S$, each with one elementary charge, pull by gravity exactly as hard as they push electrically.
  \item \emph{The Planck mass.} $m_P = \sqrt{\hbar c/G} \approx 2.17643 \times 10^{-8}$ kg, so $m_P/m_S \approx 11.706$ and $(m_S/m_P)^2 \approx 1/137.036$, the fine-structure constant (Remark~\ref{rem:stoney-planck}).
\end{itemize}
\end{workedexample}

None of these numbers is new, and honesty requires saying so. The Stoney mass follows at once from the size given to the Composer Unit, and its link with the fine-structure constant is an identity among the constants, true in any language. What is new is its place. In the cosmic mass language, Stoney's unit is the step by which every charged body climbs the Composer axis: an electron or a proton carries one such step, and a nucleus with $z$ protons carries $z$ of them.

Chapter~\ref{ch:standard-model} builds on these steps. With the angle fixed, every body becomes one triangle, with one magnitude and two angles, and all $306$ bodies of its catalogue stand on one frame: every neutral body on one line, every positive body near a second line and every negative body near a third, each tipped by its mass per charge. Thomson in 1897 and Aston in 1919 measured the ratio of charge to mass of electrons and ions \citep{Thomson1897,Aston1919}. In the cosmic mass language that ratio is an angle, and the whole of ordinary matter, from the electron to the heaviest nucleus, is one drawing. No measured mass or charge is moved to make it.

\subsection{One law for nuclei}\label{subsec:al-nuclei}

When two nuclei fuse, two old rules meet. Charge is kept exactly. Mass is not: the product is lighter than the two nuclei that made it, and the missing mass leaves as energy, $Q = \Delta m\,c^2$, as Einstein showed in 1905 \citep{Einstein1905} and as Eddington proposed in 1920 for the energy of the stars \citep{Eddington1920}. Physics states these as two rules. In the cosmic mass language they are one statement. The cosmic masses add part by part, so the charge parts of the two nuclei add exactly to the charge part of the product, and the loss of cosmic mass, $\Delta\mu = \Delta m$, is a real number: a neutral body, with no charge part at all (Proposition~\ref{prop:fusion-frame}). The charge part is never lost; only the real part is, and it leaves as the energy $Q = \Delta m\,c^2$.

Chapter~\ref{ch:closure} takes one step more. Six neighbouring nuclei, added with alternating signs, close on themselves: their charge parts cancel exactly around the loop, at every angle, and their binding nearly cancels (\Cref{thm:exact-closure}). From that one fact the mass of a nucleus can be read off its neighbours, for $2{,}481$ measured nuclei to $0.25$ MeV, and to $0.11$ MeV for those with $A \ge 40$ (Law~\ref{law:closure}), and the energy of a fusion can be known before it is weighed.

\subsection{One law for fields and motion}\label{subsec:al-fields}

Chapter~\ref{ch:unified} carries the same number into motion. Two fields whose values are General Numbers obey four equations, and one law of motion moves every body (Law~\ref{law:unified}). Opened into its two coordinates, the law is two sets of equations at once (Proposition~\ref{prop:un-split}). The Composer coordinate is Maxwell's set, with the electric and magnetic fields $\mathbf{E} = -|\Gp|^2\,\mathbf{e}$ and $\mathbf{B} = -|\Gp|^2\,\mathbf{b}$ and the magnetic constant $\mu_0 = 4\pi K/c^2$ (\Cref{thm:un-maxwell}). The real coordinate is Newton's gravity, with the magnetic part that Heaviside wrote down in 1893 \citep{Heaviside1893}. Opened in the same way, the law of motion is Newton's law of motion with gravity and the force of Lorentz in one line (Proposition~\ref{prop:un-lorentz}):
\[
m\,\frac{d\mathbf{v}}{dt} = m\,(\mathbf{g} + \mathbf{v}\times\mathbf{B}_g) + q\,(\mathbf{E} + \mathbf{v}\times\mathbf{B}) .
\]
Both parts of the field travel at the speed of light (\Cref{thm:un-waves}); one equation keeps the source, mass and charge together (\Cref{thm:un-source}); and all of it follows from a single action (\Cref{thm:un-action}). Run on Thomson's tube, the cyclotron, the drift in crossed fields, the hydrogen atom and the Moon, the law gives what Newton, Maxwell and Lorentz give (\Cref{tab:un-results}).

Here the language does something that Maxwell could not. In 1864 he saw that a field theory of gravity would give the field of gravitation a negative energy, and he could not place the sign (Chapter~\ref{ch:unified}). In the cosmic mass language one formula gives the energy of both parts (\Cref{thm:un-energy}), and opened into its coordinates it reads
\[
u = -\frac{g^2 + c^2 B_g^2}{8\pi G} + \frac{E^2 + c^2 B^2}{8\pi K}
\]
(Proposition~\ref{prop:un-two-faces}): negative for gravity, positive for the Maxwell part. The minus sign is the one that pushes like charges apart, the real part of the square of the Composer Unit. One fact of geometry gives both the sign of the force between charges and the sign of the energy of gravity. What that sign costs for the waves of gravity, Chapter~\ref{ch:unified} explains.

\begin{table}[tb]
\myfloatalign
\small
\begin{tabularx}{\textwidth}{@{}>{\raggedright\arraybackslash}X >{\raggedright\arraybackslash\hsize=1.25\hsize}X@{}}
\toprule
\tableheadline{Given} & \tableheadline{Gained} \\
\midrule
one number for each body, $\mu = m + q\,\Gp$, with its two measured parts & the opposite signs of gravity and electricity, from the plane itself \\
the shape of Newton's law, written for that number & Stoney's mass as the charge part of one elementary charge, and a picture of why electricity outweighs gravity between particles \\
the size of the Composer Unit, $|\Gp| = \sqrt{K/G}$, chosen so that the electric term has Coulomb's strength & one frame for $306$ bodies, with the charge-to-mass ratio as an angle \\
one rule for reading a product: its real part & the conservation of charge and the energy of fusion as two coordinates of one statement, and the exact closure of charge around six neighbouring nuclei \\
for moving bodies, four field equations of Maxwell's form, written for General Numbers, and one law of motion & Heaviside's gravity beside Maxwell's equations, with Maxwell's constants drawn from $G$ and $|\Gp|$, and waves at the speed of light in both \\
 & Newton's gravity and the force of Lorentz in one law of motion; the opposite signs of the two energies; mass and charge kept by one equation; all of it from one action \\
\bottomrule
\end{tabularx}
\caption[What the language was given, and what it gave back]{What the cosmic mass language was given, and what it gave back. The items on the right were not put in one by one: they come out of the few on the left.}
\label{tab:al-given}
\end{table}

\subsection{Why unification matters}\label{subsec:al-why}

Unification is one of the oldest ways in which physics has moved forward, and its history shows what kind of gain it is.

In Book~III of the \emph{Principia}, Newton compared the fall of the Moon towards the Earth with the fall of a stone at its surface, and found one force behind both \citep{Newton1687}. With one law for the heavens and the Earth, tides, comets and the shape of the planets became problems that could be calculated, and in 1846 the same law found a planet that nobody had seen (\Cref{tab:al-history}). In 1865 Maxwell joined electricity and magnetism, and found light inside them \citep{Maxwell1865}; some twenty years later Hertz made the waves that his equations described \citep{Hertz1893}, and radio followed. In 1908 Hermann Minkowski told a meeting in Cologne that space by itself and time by itself would sink into shadows, and that only a union of the two would keep its independence \citep{Minkowski1909}. On that union Einstein built his theory of gravity seven years later.

Not every unification brought a new prediction at once, and some of the most fruitful brought none. In 1788 Lagrange rewrote Newton's mechanics without a single new law, in a form built on energies instead of forces \citep{Lagrange1788}, and in 1834 Hamilton gave it another form, built on one function, the action \citep{Hamilton1834}. For the planets they predicted nothing that Newton had not. Nearly a century later Schr{\"o}dinger built wave mechanics on Hamilton's analogy between mechanics and optics \citep{Schrodinger1926}, and Feynman built a whole form of quantum mechanics on the action \citep{Feynman1948}. In 1893 Steinmetz changed no law of electricity: he wrote alternating currents with complex numbers, as in the impedance of Chapter~\ref{ch:one-number}, and the arithmetic of circuits became routine \citep{Steinmetz1893}.

Three lessons follow. A language that holds known laws in one form is shorter: fewer starting points carry more results. It explains: what looked like coincidences of form become consequences of one structure. And it is often the door through which the next theory enters, sometimes long after it was built.

\subsection{What the cosmic mass language gives}\label{subsec:al-value}

Measured against these examples, the result of this volume is of the second kind: a new form for old laws, with the values that such forms have shown before. A fair way to weigh it is to set what the language was given beside what it gave back (\Cref{tab:al-given}).


The second column is longer than the first, and its entries were not put in one by one: they come out of the few things in the first. That is the sense in which the cosmic mass language does not only keep the old laws but finds them again. For physics, five values follow.
\begin{enumerate}
  \item \emph{Economy.} Where physics uses separate laws for gravity, for electricity, for the sign of charges, for the fields and for motion, and two laws of conservation, the language uses one number for each body, one product, one rule of reading and one action.
  \item \emph{Explanation.} The shared shape of the laws of Newton and Coulomb is no longer a coincidence: they are the two parts of one product. The rule of signs, the sign of the energy of gravity and the strength of electricity over gravity become facts of one geometry.
  \item \emph{One picture of matter.} From the electron to the heaviest nucleus, every body has one card on one frame.
  \item \emph{One calculation.} One line of force serves every body, as the program of Chapter~\ref{ch:unified} shows: the kind of value that the complex impedance had for circuits.
  \item \emph{An address for the new.} Every law that has been measured lives in the real part. Anything new that the language can say must therefore live in the Composer coordinate, with the Coupling term and the angle. Faraday's hope of a link between gravity and electricity, which had no place to stand, now has an exact address.
\end{enumerate}

None of this is yet a new measurement. It is the value of a language, the kind of value that the forms of Lagrange, Hamilton and Steinmetz had before their full use was seen. Whether the cosmic mass language also carries new physics is the question of the next section, and of Volume~2.

\begin{plainwords}
Every old law this volume has tried comes out of the same few things: one number for each body, one product for each pair of bodies, and one rule for reading the product. Newton's gravity and Coulomb's law come out of the same product, with their opposite signs; so do Maxwell's equations, the force of Lorentz and the balance of mass and charge when nuclei fuse. Even an old unit of mass, Stoney's, turns up in a new place: it is the charge part of one electron. Physics has gained from such unions before, by writing less, by explaining more, and by opening doors that were found only later. That is what the cosmic mass language offers now; whether it also finds something new is the next question.
\end{plainwords}

%------------------------------------------------
\section{The Coupling Force: a candidate for new results}\label{sec:al-coupling}
%------------------------------------------------

If the language is to pay its second debt, the natural place to look is the term that neither Newton's law nor Coulomb's contains. Multiplied out, the product of two cosmic masses has three terms: mass with mass, charge with charge, and the cross term $m_a q_b + q_a m_b$, which carries the Composer Unit once (\Cref{thm:cosmic-law}). The book calls it the Coupling Force (Remark~\ref{rem:three-components}). It vanishes between neutral bodies, and between a particle and its antiparticle, which have equal masses and opposite charges.

The term is not small. Between two equal bodies it is exactly twice the geometric mean of the other two terms, so it sits halfway between gravity and electricity on a scale of powers of ten.

% keep the Coupling example on one page: break the page here if fewer than 13 lines are left
\par\vskip 0pt plus 13\baselineskip\penalty-100\vskip 0pt plus -13\baselineskip\vskip 13\baselineskip\penalty9999\vskip -13\baselineskip\vskip 0pt
\begin{workedexample}[Worked example: the size of the Coupling term]
\begin{itemize}
  \item \emph{Two protons one metre apart.} The gravity term is about $1.9 \times 10^{-64}$ N, the electric term $2.3 \times 10^{-28}$ N, and the Coupling term $4.2 \times 10^{-46}$ N $= 2\sqrt{1.9 \times 10^{-64} \times 2.3 \times 10^{-28}}$ N (\Cref{fig:two-parts}): $2\,m_S/m_p \approx 2.2 \times 10^{18}$ times gravity, and $2\,m_p/m_S \approx 1.8 \times 10^{-18}$ times the electric push.
  \item \emph{The Earth and one electron at the ground.} The gravity term is the weight of the electron, $m_e\,g \approx 8.9 \times 10^{-30}$ N. The Coupling term pairs the mass of the Earth with the charge of the electron: $g\,m_S \approx 9.8\ \mathrm{m/s^2} \times 1.859 \times 10^{-9}\ \mathrm{kg} \approx 1.8 \times 10^{-8}$ N, or $m_S/m_e \approx 2 \times 10^{21}$ times the weight.
\end{itemize}
\end{workedexample}

Measurement already says what the term is not. If any sizeable share of it acted along the line between two bodies, even the share $\cos\pb \approx 0.176$ that the reading by the shadow keeps (\Cref{sec:reading-law}), the neutron would carry an apparent charge about two hundred standard errors beyond what is measured (\Cref{subsec:test-sign}), and every free electron near the ground would feel a force more than $10^{20}$ times its weight, which any electron beam in a laboratory would show. Read by its real part, the law gives the term no part along the line, which is what these measurements require. Whatever the Coupling term means, it is not a force between the mass of one body and the charge of another of anything like its own size.

Where, then, could it show itself? Read by its real part, the Coupling term lives in the Composer coordinate of the product, together with the angle of the Composer axis (Definition~\ref{def:reading}). The same holds for bodies in motion: the law of motion of Chapter~\ref{ch:unified} reads only the real part of $\mu\,(\boldsymbol{\mathcal{E}} + \mathbf{v}\times\boldsymbol{\mathcal{B}})$, and its Composer coordinate, which pairs the mass of the body with the charge part of the field and its charge with the gravity part, is left unread (Law~\ref{law:unified}).

Volume~2 begins from this. Its question is whether any measurable quantity reads the Composer coordinate of a product. For slow bodies, Chapter~\ref{ch:unified} shows that forces, energies and momenta do not. What remains are the places this volume does not reach: fast bodies and strong fields, where its law gives way to general relativity, and the quantum world, where what is measured is not always a force. A prediction found there must meet the standard of \Cref{sec:al-duties}:
\begin{enumerate}
  \item a number, with its error, drawn from the law with no new constant fitted to the effect it predicts;
  \item the conditions under which it holds, and an experiment able to see it;
  \item agreement with every measurement that already limits the term, among them those of the neutron and of hydrogen;
  \item the result that would refute it, stated before the measurement is made.
\end{enumerate}
No such prediction is claimed in this volume. The Coupling term is a candidate, not a result.

The second open question meets the first in the same place. No force fixes the angle of the Composer axis, because the real part of a product does not contain it (\Cref{subsec:which-root}); the angle lives, with the Coupling term, in the Composer coordinate. A measurement that reads that coordinate, if one exists, would be the first that could fix the angle, and with it the choice of $79.86^\circ$.

One set of predictions in this volume is already at risk, and it should be named. Chapter~\ref{ch:closure} registers the masses of $61$ nuclei that have not been measured, each with its stated error, after a first list of $67$ in which $24$ masses turned out to be measured already, and $20$ of those fell within one stated error (\Cref{sec:cl-predictions}). These predictions test the Closure Law, and with it the claim that mass and charge are kept as the two coordinates of one number. They do not test the Coupling term, or the angle.

\begin{plainwords}
The Coupling term pairs the mass of one body with the charge of the other. As a number it is large: between the Earth and an electron, far larger than the electron's weight. As a force it cannot exist, or laboratories would have seen it long ago, and the real reading gives it no force at all. It lives instead in the second part of the product, together with the angle. Volume~2 asks whether anything we can measure reads that second part. If something does, the Coupling term could give the first new prediction of the theory; if nothing does, the cosmic mass remains a language for the physics we already know.
\end{plainwords}

%------------------------------------------------
\section*{Two volumes}\phantomsection\label{sec:al-close}
\addcontentsline{toc}{section}{Two volumes}
\markright{\spacedlowsmallcaps{Two volumes}}
%------------------------------------------------

The Coda of this book lists what every new kind of number has needed: a question that the old numbers could not answer, a definition that contradicts nothing, consequences that nobody put in on purpose, and in the end a fact about the world that could not be said without them. This volume offers a candidate for each of the first three. Its question is the one of Chapter~\ref{ch:one-number}: can one number hold the mass and the charge of a particle? Its definition is the General Number, which contains the complex numbers as one of its planes. Its consequences include the minus sign that pushes like charges apart, which comes from the plane itself, and the laws gathered in \Cref{tab:al-found}, which the language gives back in one form.

The fourth need is the work of Volume~2. There the Coupling term and the angle of the Composer axis will be asked for numbers that can be measured, and those numbers will be set against the world.

The rest of this volume builds the language and tests it. Chapters~\ref{ch:oblique} and~\ref{ch:general-numbers} build the numbers, from a triangle that leans to a plane whose second axis leans. Chapter~\ref{ch:prologue} explains why the language of physics must sometimes grow. Chapter~\ref{ch:interaction} writes the law for two bodies and the rule for reading it, Chapter~\ref{ch:standard-model} draws every particle, ion and nucleus on one frame, Chapter~\ref{ch:closure} weighs the nuclei, and Chapter~\ref{ch:unified} writes one law for fields and moving bodies. The next chapter begins with the oldest piece of that road: a triangle that does not stand upright.

\begin{plainwords}[What we have so far]
\begin{itemize}
  \item A new theory must first give back the old laws, exactly, wherever they have been tested, and then predict something new that a measurement could prove wrong.
  \item This volume pays the first debt: one number for each body, read by its real part, gives back the laws of Newton, Coulomb, Maxwell and Lorentz, Heaviside's gravity and the balance of nuclear fusion, and it gives Stoney's unit of mass a new place.
  \item Such a union has value before any new prediction: fewer starting points, coincidences explained, one picture of matter, one calculation for every body, and an exact address for what is new.
  \item None of this depends on the angle of the Composer axis, and none of it is a new prediction.
  \item The candidate for the second debt is the Coupling term, which lives with the angle in the Composer coordinate of a product. Finding what, if anything, measures that coordinate is the task of Volume~2.
\end{itemize}
\end{plainwords}
